\chapter{Calculus and Optimization}

Training a model means finding parameters that minimize a loss function.
Calculus tells us which direction to move in; optimization algorithms tell
us how to get there efficiently.

\section{Derivatives and Gradients}

\begin{definition}[Gradient]
For a scalar function $f:\mathbb{R}^n \to \mathbb{R}$, the gradient is the
vector of partial derivatives
\[
  \nabla f(\mathbf{x}) = \left(\frac{\partial f}{\partial x_1}, \ldots, \frac{\partial f}{\partial x_n}\right).
\]
\end{definition}

The gradient points in the direction of \emph{steepest ascent} of $f$. To
minimize a loss function, we move in the direction of $-\nabla f$.

\section{The Chain Rule}

\begin{definition}[Chain Rule]
If $z = f(y)$ and $y = g(x)$, then
\[
  \frac{dz}{dx} = \frac{dz}{dy}\cdot\frac{dy}{dx}.
\]
\end{definition}

The chain rule is the mathematical backbone of backpropagation: a neural
network is a composition of many functions (layers), and computing the
gradient of the loss with respect to early-layer weights requires
multiplying together the derivatives of every layer in between.

\section{Gradient Descent}

Given a loss function $J(\theta)$, gradient descent iteratively updates
parameters:
\[
  \theta^{(t+1)} = \theta^{(t)} - \eta\,\nabla J(\theta^{(t)})
\]
where $\eta$ is the \textbf{learning rate}, controlling step size.

\figw{gradient_descent_contour.png}{0.68}{Gradient descent path on a convex bowl-shaped loss $J(\theta_1,\theta_2)=\theta_1^2+3\theta_2^2$. Each step moves opposite the gradient, converging toward the minimum.}

\figw{loss_surface_3d.png}{0.75}{The same loss function viewed as a 3D surface, with the descent trajectory rolling downhill toward the minimum.}

\begin{keyidea}
The learning rate $\eta$ trades off speed and stability. Too small and
training crawls; too large and updates overshoot the minimum or diverge
entirely. Elongated, narrow bowls (as in the figure above, where one
direction is steeper than the other) are exactly why adaptive optimizers
like Adam and RMSProp were developed — they rescale step sizes per
parameter.
\end{keyidea}

\section{Variants of Gradient Descent}

\begin{itemize}[leftmargin=*]
  \item \textbf{Batch GD:} uses the entire dataset per update — accurate but slow.
  \item \textbf{Stochastic GD (SGD):} uses one example per update — noisy but fast, and the noise can help escape shallow local minima.
  \item \textbf{Mini-batch GD:} the standard compromise, updating on small batches (e.g.\ 32–256 examples).
  \item \textbf{Momentum:} accumulates a moving average of past gradients to smooth out the path and accelerate convergence.
  \item \textbf{Adam:} combines momentum with a per-parameter adaptive learning rate, and is the default optimizer in most modern deep learning.
\end{itemize}

\section{Convexity}

\begin{definition}[Convex Function]
$f$ is convex if for all $\mathbf{x}, \mathbf{y}$ and $t \in [0,1]$,
\[
  f(t\mathbf{x} + (1-t)\mathbf{y}) \le t f(\mathbf{x}) + (1-t) f(\mathbf{y}).
\]
\end{definition}

For convex loss functions (like linear regression's MSE), gradient descent
is guaranteed to converge to the \emph{global} minimum. Neural network loss
surfaces are generally \emph{non}-convex, which is why training deep
networks well requires careful initialization, learning-rate schedules, and
architectural choices — there's no guarantee of finding the global optimum.
